\section{Action, gravity, and thermal information}
\label{sec:gravity}
\subsection{Dodecaphasic period and circular realization of action}
The calendar already constructed determines $T_{108}=108t_0$.
In its circular realization,
\[
\omega_{108}=\frac{2\pi}{108t_0},\qquad
R_{108}=\frac{c_{\rm int}}{\omega_{108}}
=\frac{54}{\pi}\ell_0,\qquad \ell_0=c_{\rm int}t_0.
\]
The unit $\ell_0$ and the time and velocity charts remain
explicit; the SI chart does not select the cycle.
The longitudinal composition in Section~\ref{md:rigidez:longitud}
fixes $R_{108}=L_\alpha$ and
$t_0=\pi_{\rm HMT}L_\alpha/(54c_{\rm int})$ in this edition.
The following circular argument is retained as a specialization
of the unit character in Theorem~\ref{md:rigidez:teorema}.
For each oriented action section $a\in\{\mathrm{pre},\mathrm{ret}\}$,
\[
E_{108,a}=\hbar_a\omega_{108},\qquad
m_{108,a}=\frac{E_{108,a}}{c_{\rm int}^2}.
\]
The reduced and full lengths are, respectively,
\[
\frac{\hbar_a}{m_{108,a}c_{\rm int}}=R_{108},
\qquad
\frac{h_a}{m_{108,a}c_{\rm int}}=2\pi R_{108}.
\]
The same action and duration enter both readings.

\begin{proposition}[Radius-coincidence selector]
In the dimensional collection
\[
G_a(\vartheta)=\vartheta\frac{c_{\rm int}^5t_0^2}{\hbar_a},
\qquad \vartheta>0,
\]
the circular-realization condition
$G_a(\vartheta)m_{108,a}/c_{\rm int}^2=R_{108}$
has the unique solution
\[
\vartheta^\circ_{108}=(54/\pi)^2.
\]
\end{proposition}
\begin{proof}
The first radius is
$r_g(\vartheta)=\vartheta(\pi/54)\ell_0$.
Comparing it with $(54/\pi)\ell_0$, with $\ell_0>0$, reduces the
equality to a strictly increasing linear equation.
\end{proof}
The coincidence condition is an explicit part of this realization;
the theorem solves for its coefficient without using a metrological value of $G$.
The radius used here is $Gm/c^2$, so the factor two from
the convention $2Gm/c^2$ is not inadvertently transferred.

\begin{proposition}[Reciprocal action and gravity, preserved area]
In the two preceding charts,
\[
\frac{G_{\rm pre}}{G_{\rm ret}}=R_{\rm act}^{-1},\qquad
\hbar_aG_a=\vartheta^\circ_{108}c_{\rm int}^5t_0^2,
\qquad
\ell_P^2=\frac{\hbar_aG_a}{c_{\rm int}^3}
=\vartheta^\circ_{108}\ell_0^2.
\]
\end{proposition}
\begin{proof}
The definition of $G_a$ contains $\hbar_a^{-1}$, and the remaining factors
are common to both charts. Taking the quotient and product gives
the three equalities.
\end{proof}
